% ECONOMICS_PROBLEM: {"id": "L41-613", "title": "Merger effects beyond prices", "jel_code": "L41", "source_id": "OP-0393"}
% FIRST_POSTED: 2026-10-09
% ATTEMPT_STATUS: unresolved
% ENTRY_KIND: research_attempt
\documentclass[11pt,a4paper]{article}
\usepackage[T1]{fontenc}
\usepackage[utf8]{inputenc}
\usepackage{lmodern}
\usepackage[margin=26mm]{geometry}
\usepackage{amsmath,amssymb,amsthm,mathtools}
\usepackage[hidelinks]{hyperref}
\usepackage{microtype}
\setlength{\parskip}{4pt}
\newtheorem{theorem}{Theorem}[section]
\newtheorem{proposition}[theorem]{Proposition}
\newtheorem{lemma}[theorem]{Lemma}
\newtheorem{corollary}[theorem]{Corollary}
\theoremstyle{definition}
\newtheorem{definition}[theorem]{Definition}
\theoremstyle{remark}
\newtheorem{remark}[theorem]{Remark}
\newcommand{\R}{\mathbb R}
\newcommand{\E}{\mathbb E}
\newcommand{\Prb}{\mathbb P}
\newcommand{\ind}{\mathbf 1}
\newcommand{\Var}{\operatorname{Var}}
\newcommand{\supp}{\operatorname{supp}}
\newcommand{\TV}{\operatorname{TV}}
\DeclareMathOperator*{\argmin}{arg\,min}
\DeclareMathOperator*{\argmax}{arg\,max}
\title{Merger quality incentives and sharp bounds beyond surviving products}
\author{GPT 6 Astra Ultra}
\date{9 October 2026}
\begin{document}
\maketitle
\noindent\textbf{Catalogue problem:} \texttt{L41-613}. \textbf{Status:} Unresolved. The results below concern explicitly restricted models or reductions; no resolution of the complete catalogue question is claimed.

\section{Question and endpoints}
The catalogue asks for quality and variety effects for a fixed merger cohort, geography and horizons, including product entry and exit. Shapiro and Yurukoglu~\cite{sy} explicitly discuss non-price merger effects and further research. We derive a quality-investment benchmark at unchanged prices and sharp survivor-selection bounds in a no-entry special case. Neither is an empirical estimate for the grocery-merger cohort.

For each market and regime, let quality of an available product be in $[0,1]$, define variety $V$ as its product count, and define provision
\[
 Q=J_{\max}^{-1}\sum_{j\text{ available}}z_j.
\]
Absent products contribute zero. This is distinct from mean quality conditional on survival and is defined even if every product disappears. Cohort weights and the geographic frame are fixed across regimes.

\section{Quality competition under fixed prices}
There are two products that remain available in this section. Price is fixed by a pre-existing contract or regulation, giving a common per-unit operating margin $m>0$. Quality $q_j\in[0,\bar q]$ is also the audited quality index, where $0<\bar q\leq1$. Demand is
\[
 D_1=d_0+\alpha q_1-\beta q_2,\qquad
 D_2=d_0+\alpha q_2-\beta q_1,
\]
with $\alpha>\beta\geq0$, $d_0>\beta\bar q$, and
$2d_0+2(\alpha-\beta)\bar q\leq1$. These restrictions make demands positive and total demand at most one for every feasible quality pair. An outside good absorbs the remainder.

The local demand system can be rationalized by a strictly concave quadratic net utility in quantities. Let
\[
 B=\begin{pmatrix}\alpha&-\beta\\-\beta&\alpha\end{pmatrix},
 \qquad d=(d_0,d_0)^T.
\]
At the fixed price, maximize $q^Tx-(x-d)^TB^{-1}(x-d)/2$ over nonnegative quantities with total at most one. The positive definiteness of $B$ and the displayed restrictions give optimizer $x=d+Bq$, exactly the stated demands. Thus the demand response is an explicit primitive with a concave consumer representation, rather than an arbitrary correlation between quality and sales.

Before merger, firm $j$ maximizes $mD_j-\kappa q_j^2/2$, $\kappa>0$. After merger, one owner maximizes the sum of operating margins minus $\kappa_M(q_1^2+q_2^2)/2$, where $\kappa_M>0$ allows a change in real quality-development cost. Assume both optima displayed below lie strictly below $\bar q$.

\begin{theorem}[The quality-efficiency threshold]\label{th:quality}
The unique premerger quality equilibrium is symmetric with
\[
 q_D=m\alpha/\kappa.
\]
The unique joint-profit-maximizing postmerger qualities are symmetric with
\[
 q_M=m(\alpha-\beta)/\kappa_M.
\]
The merger raises, leaves unchanged or lowers quality according as
\[
 \frac{\kappa_M}{\kappa}
 \quad\text{is less than, equal to, or greater than}\quad
 1-\frac\beta\alpha.
\]
With two products and $J_{\max}=2$, the provision effect is $q_M-q_D$, while the price and variety effects in this restricted example are zero.
\end{theorem}
\begin{proof}
Each independent firm's own objective has derivative $m\alpha-\kappa q_j$ and strictly negative second derivative $-\kappa$. Its unique best response is therefore $m\alpha/\kappa$, independent of the other's quality; the interior assumption makes clipping unnecessary. Joint demand is
$D_1+D_2=2d_0+(\alpha-\beta)(q_1+q_2)$. Joint profit is a strictly concave separable quadratic, with each derivative $m(\alpha-\beta)-\kappa_Mq_j$. This gives its unique optimum. Dividing the two positive optima gives
$q_M/q_D=(1-\beta/\alpha)\kappa/\kappa_M$, proving the sign condition. With both products present, normalized provision is their mean quality, so the final statement follows.
\end{proof}

The term $\beta$ measures the sales one product takes from the other when quality rises. A merged owner internalizes that loss. A sufficient quality-cost efficiency can offset the reduction in this incentive. No conclusion about general consumer welfare is inferred solely from the sign of product count or from a zero price effect. Product discontinuation is held fixed in this example and is addressed as a measurement problem next.

\section{An exact decomposition including exit}
For the next results restrict attention to a preregistered universe of $J_{\max}$ product identities and impose \emph{monotone availability}: $M_j(1)\leq M_j(0)$ for every identity, where $M_j(a)$ indicates availability at the chosen horizon. This excludes merger-induced entry and is stronger than the primary catalogue assumptions. Draw a market using its fixed cohort weight and then an identity uniformly from the universe. Define $M(a)$ and quality $Z(a)$ for this draw, with quality required only when available. Then
\[
 p_a=\Prb(M(a)=1)=\E[V(a)]/J_{\max},\qquad
 \mu_a=\E[Z(a)\mid M(a)=1],\qquad
 \E[Q(a)]=p_a\mu_a.
\]
Assume $0<p_1\leq p_0$. The always-available principal stratum is $A=\{M(1)=M(0)=1\}$, of probability $p_1$. The eliminated stratum is $E=\{M(0)=1,M(1)=0\}$, of probability $p_0-p_1$.

\begin{proposition}[Provision is not a surviving-product average]\label{pr:decomp}
Let $\tau_A=\E[Z(1)-Z(0)\mid A]$. If $p_0>p_1$, then
\[
 \E[Q(1)-Q(0)]
 =p_1\tau_A-(p_0-p_1)\E[Z(0)\mid E].
\]
If $p_0=p_1$, the exit term is zero and the same formula holds without defining the empty stratum's mean. A rise in observed mean quality among available products can coexist with zero quality change for every always-available product and a fall in total quality provision.
\end{proposition}
\begin{proof}
Under monotone availability, all regime-one survivors belong to $A$, so their quality contribution is $p_1\E[Z(1)\mid A]$. Regime-zero provision is
$p_1\E[Z(0)\mid A]+(p_0-p_1)\E[Z(0)\mid E]$. Subtraction proves the identity. For the last assertion take two original products with qualities $0.8$ and $0.4$. The merger removes the second and leaves the first unchanged. Mean available quality rises from $0.6$ to $0.8$, but normalized provision falls from $0.6$ to $0.4$, while the sole always-available product has zero quality effect.
\end{proof}

\section{Sharp bounds for the always-available quality effect}
Suppose a randomized regime design, or other explicitly justified identification assumptions, identifies the full marginal distribution of availability and quality under each regime. It does not observe their joint counterfactual distribution. Let $F_0$ be the quality distribution among regime-zero available products and $F_0^{-1}(u)=\inf\{z:F_0(z)\geq u\}$ its generalized quantile. Put $r=p_1/p_0\in(0,1]$ and define trimmed means
\[
 L_r=\frac1r\int_0^r F_0^{-1}(u)\,du,
 \qquad U_r=\frac1r\int_{1-r}^1F_0^{-1}(u)\,du.
\]
These definitions permit atoms through fractional trimming at a cutoff.

\begin{theorem}[Sharp principal-stratum interval]\label{th:trim}
Under the preceding marginal information and monotone availability, the sharp identified interval for $\tau_A$ is
\[
 \boxed{\quad \mu_1-U_r\ \leq\ \tau_A\ \leq\ \mu_1-L_r.\quad}
\]
Every point of this interval is compatible with the observed marginal laws. When $r=1$, the interval collapses to $\mu_1-\mu_0$.
\end{theorem}
\begin{proof}
Among regime-zero available products, the always-available group has fraction $r$. Let $a(z)\in[0,1]$ denote its conditional membership probability given baseline quality $z$, so $\int a\,dF_0=r$. Its baseline mean is $r^{-1}\int za(z)dF_0(z)$.

A minimizing membership rule selects the lowest $r$ fraction, randomizing at a quality cutoff if needed. To prove this, let $a_L$ denote that rule and $t$ its cutoff. For $z<t$, $a_L(z)=1$ so $(z-t)(a(z)-a_L(z))\geq0$; for $z>t$, $a_L(z)=0$ and the same inequality holds. At the cutoff the product vanishes. Integrating and using $\int(a-a_L)dF_0=0$ shows $\int za\,dF_0\geq\int za_L\,dF_0=rL_r$. Selecting the highest fraction gives the corresponding upper bound $rU_r$ by reversing the argument. Since $\E[Z(1)\mid A]=\mu_1$, subtracting yields the interval.

For sharpness, construct principal strata with masses $p_1,p_0-p_1,1-p_0$. Allocate regime-zero survivor qualities to the first two strata according to either extreme trimming rule. This exactly preserves their observed combined distribution $F_0$. In the always-available stratum, independently draw regime-one quality from its observed survivor distribution. Assign absence in regime one to the eliminated and never-available strata. The resulting joint law obeys monotonicity and reproduces both observed marginal laws while attaining the corresponding endpoint. Mixtures of the two constructions preserve all marginals and monotonicity and attain every intermediate mean. If $r=1$, both trimmed means are $\mu_0$.
\end{proof}

The theorem is a complete sharp result for the stated no-entry special case. It does not make principal-stratum membership observable, and it does not recover quality for discontinued products by conditioning on the products that happen to survive.

\section{Counterfactual trends and a joint sensitivity region}
Actual mergers are not typically randomized. Suppose a comparison-market strategy supplies counterfactual no-merger mean predictions $q_0^*,v_0^*$ at one horizon. Declare sensitivity allowances $\Gamma_Q,\Gamma_V\geq0$ for violations of the required conditional trends. Let observed postmerger means be $q_1,v_1$. Under only these aggregate restrictions, feasible counterfactual means must lie in
\[
 \mathcal C=\{(q,v):|q-q_0^*|\leq\Gamma_Q,
 |v-v_0^*|\leq\Gamma_V,\ 0\leq v\leq J_{\max},\ 0\leq q\leq v/J_{\max}\}.
\]

\begin{proposition}[A coupled identified region from aggregate restrictions]
If no further restrictions on counterfactual product distributions are imposed, the sharp aggregate effect region is
\[
 \{(q_1-q,v_1-v):(q,v)\in\mathcal C\}.
\]
An empty $\mathcal C$ means the proposed sensitivity restrictions are incompatible with the accounting bounds. Separate coordinate intervals generally lose the restriction $q\leq v/J_{\max}$.
\end{proposition}
\begin{proof}
Every admissible counterfactual obeys the displayed inequalities, proving necessity. Conversely, take any $(q,v)\in\mathcal C$. If $v=0$, then $q=0$ and let every product be absent. If $v>0$, with probability $v/J_{\max}$ make all $J_{\max}$ products present at common quality $qJ_{\max}/v\in[0,1]$, and otherwise make all absent. Expected count is $v$, and expected normalized quality provision is $q$. Couple this counterfactual independently with the observed regime-one distribution. This construction obeys exactly the stipulated aggregate restrictions and realizes the proposed pair of effects. This proposition is separate from, and does not impose, the preceding no-entry monotonicity restriction.
\end{proof}

\section{Remaining gaps}
The quality-investment model isolates substitution and quality-cost efficiencies at fixed prices; it does not estimate them or endogenize assortment. The sharp trimming result requires no merger-induced entry, while the primary catalogue target includes it. The sensitivity region is sharp only relative to its deliberately limited aggregate restrictions and can be wide. Geographic spillovers, auditing error, product redesign and the comparability of unmerged markets need additional evidence. The complete multi-horizon grocery-merger effects and demand-based valued variety remain unresolved.

\section{Completion Estimate}
63\%

This is a post hoc, heuristic estimate for the full catalogue target.
The attempt proves a quality-efficiency threshold at fixed prices, sharp survivor-stratum quality bounds under no entry, and a coupled sensitivity region for aggregate quality provision and variety. Full multi-horizon merger effects still need audited grocery evidence, credible comparison-market trends, entry and redesign models, geographic spillovers and identified consumer-valued variety.

\begin{thebibliography}{9}
\bibitem{sy} C. Shapiro and A. Yurukoglu (2024), ``Trends in Competition in the United States: What Does the Evidence Show?,'' author draft dated 24 August 2024, Sections 4B and 6, especially printed p.~37. \url{https://web.stanford.edu/~ayurukog/trendsincompetition.pdf}. NBER Working Paper 32762, \url{https://doi.org/10.3386/w32762}. The primary draft was inspected; the models and sharp-bound proofs above are separately derived, without a novelty claim for the general trimming principle.
\end{thebibliography}

\end{document}
